% GENERATED FILE. Edit canonical lessons, then run npm run generate:books.
% canonical-source-hash: fnv1a64:7e5ee1a8c867b63b
% canonical-lessons: TE-C125-pavuram, TE-C125-gudlaguba, TE-C125-nemali, TE-C125-sitakokacilaka, TE-C125-kaki

\chapter{Pigeon, Owl, Peacock, Butterfly, Crow}
\label{ch:te-pigeon-owl-peacock-butterfly-crow}
\hlchaptermodality{125}

\begin{chapteropening}
\textbf{By the end of this chapter:} \emph{I can say a pigeon, an owl, a peacock, a butterfly and a crow.}

Complete the last lesson of chapter 125: I can say a pigeon, an owl, a peacock, a butterfly and a crow.
\end{chapteropening}

\section[\texorpdfstring{\te{పావురం} (pāvuraṁ)}{పావురం (pāvuraṁ)}]{\te{పావురం} (pāvuraṁ) — a pigeon}

\label{lesson:TE-C125-pavuram}



\begin{quote}
Before the new one: say the Telugu for a tortoise, then the Telugu for a parrot.
\end{quote}

\subsection*{You'll want to know: \te{పావురం}}
\textbf{\te{పావురం}} — \emph{pāvuraṁ} — \textquotedblleft{}a pigeon\textquotedblright{}.

\begin{grammarlens}[title={What you've built}]
One more word for this chapter.
\end{grammarlens}

\subsection*{Your turn}
\begin{itemize}
\raggedright
  \item \emph{Say it:} \emph{pāvuraṁ}
  \item \emph{Say it:} \emph{pāvuraṁ}, once more
  \item \emph{Say it:} say \emph{pāvuraṁ}
  \item \emph{Recall:} say the Telugu for a tortoise, then the Telugu for a parrot, then say \emph{pāvuraṁ} again
\end{itemize}

\begin{tcolorbox}[breakable,colback=teal!4,colframe=teal!35!black,title={Before you move on}]
What does \te{పావురం} mean? (\textquotedblleft{}A pigeon\textquotedblright{}.) Say it once more.
\end{tcolorbox}
\section[\texorpdfstring{\te{గుడ్లగూబ} (guḍlagūba)}{గుడ్లగూబ (guḍlagūba)}]{\te{గుడ్లగూబ} (guḍlagūba) — an owl}

\label{lesson:TE-C125-gudlaguba}



\begin{quote}
Before the new one: say the Telugu for a parrot, then the Telugu for a pigeon.
\end{quote}

\subsection*{You'll want to know: \te{గుడ్లగూబ}}
\textbf{\te{గుడ్లగూబ}} — \emph{guḍlagūba} — \textquotedblleft{}an owl\textquotedblright{}.

\begin{grammarlens}[title={What you've built}]
One more word for this chapter.
\end{grammarlens}

\subsection*{Your turn}
\begin{itemize}
\raggedright
  \item \emph{Say it:} \emph{guḍlagūba}
  \item \emph{Say it:} \emph{guḍlagūba}, once more
  \item \emph{Say it:} say \emph{guḍlagūba}
  \item \emph{Recall:} say the Telugu for a parrot, then the Telugu for a pigeon, then say \emph{guḍlagūba} again
\end{itemize}

\begin{tcolorbox}[breakable,colback=teal!4,colframe=teal!35!black,title={Before you move on}]
What does \te{గుడ్లగూబ} mean? (\textquotedblleft{}An owl\textquotedblright{}.) Say it once more.
\end{tcolorbox}
\section[\texorpdfstring{\te{నెమలి} (nemali)}{నెమలి (nemali)}]{\te{నెమలి} (nemali) — a peacock}

\label{lesson:TE-C125-nemali}



\begin{quote}
Before the new one: say the Telugu for a pigeon, then the Telugu for an owl.
\end{quote}

\subsection*{You'll want to know: \te{నెమలి}}
\textbf{\te{నెమలి}} — \emph{nemali} — \textquotedblleft{}a peacock\textquotedblright{}.

\begin{grammarlens}[title={What you've built}]
One more word for this chapter.
\end{grammarlens}

\subsection*{Your turn}
\begin{itemize}
\raggedright
  \item \emph{Say it:} \emph{nemali}
  \item \emph{Say it:} \emph{nemali}, once more
  \item \emph{Say it:} say \emph{nemali}
  \item \emph{Recall:} say the Telugu for a pigeon, then the Telugu for an owl, then say \emph{nemali} again
\end{itemize}

\begin{tcolorbox}[breakable,colback=teal!4,colframe=teal!35!black,title={Before you move on}]
What does \te{నెమలి} mean? (\textquotedblleft{}A peacock\textquotedblright{}.) Say it once more.
\end{tcolorbox}
\section[\texorpdfstring{\te{సీతాకోకచిలుక} (sītākōkacilaka)}{సీతాకోకచిలుక (sītākōkacilaka)}]{\te{సీతాకోకచిలుక} (sītākōkacilaka) — a butterfly}

\label{lesson:TE-C125-sitakokacilaka}



\begin{quote}
Before the new one: say the Telugu for an owl, then the Telugu for a peacock.
\end{quote}

\subsection*{You'll want to know: \te{సీతాకోకచిలుక}}
\textbf{\te{సీతాకోకచిలుక}} — \emph{sītākōkacilaka} — \textquotedblleft{}a butterfly\textquotedblright{}.

\begin{grammarlens}[title={What you've built}]
One more word for this chapter.
\end{grammarlens}

\subsection*{Your turn}
\begin{itemize}
\raggedright
  \item \emph{Say it:} \emph{sītākōkacilaka}
  \item \emph{Say it:} \emph{sītākōkacilaka}, once more
  \item \emph{Say it:} say \emph{sītākōkacilaka}
  \item \emph{Recall:} say the Telugu for an owl, then the Telugu for a peacock, then say \emph{sītākōkacilaka} again
\end{itemize}

\begin{tcolorbox}[breakable,colback=teal!4,colframe=teal!35!black,title={Before you move on}]
What does \te{సీతాకోకచిలుక} mean? (\textquotedblleft{}A butterfly\textquotedblright{}.) Say it once more.
\end{tcolorbox}
\section[\texorpdfstring{\te{కాకి} (kāki)}{కాకి (kāki)}]{\te{కాకి} (kāki) — a crow}

\label{lesson:TE-C125-kaki}



\begin{quote}
Before the new one: say the Telugu for a peacock, then the Telugu for a butterfly.
\end{quote}

\subsection*{You'll want to know: \te{కాకి}}
\textbf{\te{కాకి}} — \emph{kāki} — \textquotedblleft{}a crow\textquotedblright{}.

\begin{grammarlens}[title={What you've built}]
One more word for this chapter.
\end{grammarlens}

\subsection*{Your turn}
\begin{itemize}
\raggedright
  \item \emph{Say it:} \emph{kāki}
  \item \emph{Say it:} \emph{kāki}, once more
  \item \emph{Say it:} say \emph{kāki}
  \item \emph{Recall:} say the Telugu for a peacock, then the Telugu for a butterfly, then say \emph{kāki} again
\end{itemize}

\begin{tcolorbox}[breakable,colback=teal!4,colframe=teal!35!black,title={Before you move on}]
What does \te{కాకి} mean? (\textquotedblleft{}A crow\textquotedblright{}.) Say it once more.
\end{tcolorbox}
